\section{Верификация и емпирична оценка}
\label{sec:verification_evaluation}

Тази глава обединява двата вида доказателства, върху които се опира настоящата работа. От една страна стоят compile-time ограниченията, unit тестовете, integration тестовете и live backend проверките, чрез които се валидира correctness на архитектурата. От друга страна стои емпиричната оценка на асинхронния слой, чиято цел е да покаже кога coroutine-based изпълнението носи реална практическа полза и кога неговата цена остава сравнима или по-висока от синхронната алтернатива.

Важният методологичен принцип е, че тези два вида доказателства не се разглеждат като взаимозаменяеми. Тестовата верификация отговаря на въпроса дали библиотеката е коректна спрямо собствените си contracts и backend ограничения. Benchmark оценката отговаря на различен въпрос: носи ли асинхронният execution модел измерима полза при натоварвания, в които waiting фазата доминира над локалната изчислителна работа. Само съвместното присъствие на двата пласта прави архитектурните претенции на разработката достатъчно убедителни.

\subsection{Стратегия на верификацията}

Верификационната стратегия следва самата архитектура на библиотеката. Compile-time слоят се проверява чрез concept-и, capability ограничения и типизирано изграждане на query IR. Runtime слоят се проверява чрез unit тестове за рендиране, параметризиране, thread safety, миграции, корутинна инфраструктура и асинхронен достъп. Най-външният слой се валидира чрез integration тестове с реални backend-и, така че да не се смесват mock коректност и реално драйверно поведение.

В текущата ревизия на хранилището пълният автоматизиран suite включва 271 теста при изпълнение чрез \code{ctest}. Това число само по себе си не е научен резултат, но е важно като индикатор, че изложените в предходните глави механизми не са оставени на ниво концептуална схема. Те са покрити от систематични проверки по слоеве и по backend семейства.

\begin{table}[H]
\centering
\small
\caption{Основни източници на верификационни доказателства}
\label{tab:evaluation_evidence_matrix}
\begin{tabularx}{\textwidth}{L{3.1cm}L{5.0cm}Z}
\toprule
\textbf{Източник} & \textbf{Какво валидира} & \textbf{Представителни артефакти} \\
\midrule
Compile-time ограничения и query IR & типова коректност на заявки, capability-базирано ограничаване и connector contracts & \path{tests/unit/test_query.cpp}, \path{tests/unit/test_postgresql_connector.cpp}, \path{tests/unit/test_mongodb_connector.cpp} \\
Асинхронна инфраструктура & \code{Task<T>}, \code{ThreadPool}, \code{IoContext}, \code{async\_db<DB>}, coroutine-aware pooling и транзакции & \path{tests/unit/test_task.cpp}, \path{tests/unit/test_thread_pool.cpp}, \path{tests/unit/test_io_context.cpp}, \path{tests/unit/test_async_connector.cpp} \\
Миграции и thread safety & schema diff логика, rollback semantics, connection guards и конкурентна безопасност & \path{tests/unit/test_schema_migration.cpp}, \path{tests/unit/test_thread_safety.cpp} \\
Live backend интеграции & поведение върху реални SQLite, PostgreSQL, MySQL, MongoDB, Redis, Neo4j и Cassandra среди & \path{tests/integration/test_sqlite.cpp}, \path{tests/integration/test_postgresql_live.cpp}, \path{tests/integration/test_mongodb_live.cpp}, \path{tests/integration/test_redis_live.cpp} \\
Емпирична оценка & throughput при симулирана waiting фаза и изолирани coroutine overhead измервания & \path{benchmarks/bench_async.cpp} \\
\bottomrule
\end{tabularx}
\end{table}

Таблица~\ref{tab:evaluation_evidence_matrix} показва, че доказателствената база е хетерогенна по необходимост. Една compile-time ORM библиотека не може да бъде защитена само с throughput числа, както и асинхронна архитектура не може да бъде защитена само с concept проверки. Изискват се и двата вида доказателства.

\subsection{Каталог на unit и integration доказателствата}

За да не остане verification стратегията само на равнище обща рамка, тя трябва да бъде разгърната и като конкретен каталог от репозиторийни артефакти. Това е особено важно за настоящата работа, защото архитектурните претенции са distributed между entity слоя, query IR-а, connector contract-а, async инфраструктурата и live backend проверките. Следващите таблици правят именно това: показват къде в хранилището са разположени основните доказателства и каква роля играе всяка тестова група.

\begin{table}[H]
\centering
\small
\caption{Основни unit тестове за generic, core и async подсистемите}
\label{tab:verification_unit_core}
\begin{tabularx}{\textwidth}{L{3.2cm}L{3.8cm}Y}
\toprule
\textbf{Файл} & \textbf{Подсистема} & \textbf{Какво валидира} \\
\midrule
\path{test_types.cpp} & базови типови alias-и & коректност на публичните alias-и и базови utility зависимости \\
\path{test_property.cpp} & entity property layer & compile-time модел на скаларните полета и type-trait поведението \\
\path{test_relationship.cpp} & entity relationship layer & infer логика за връзките и storage-strategy поведението \\
\path{test_query.cpp} & query IR и fluent API & коректност на \code{select}, \code{insert}, \code{update}, \code{deleteq} и производните query структури \\
\path{test_schema_migration.cpp} & migration layer & DDL generation, diff логика и state-machine поведение \\
\path{test_thread_safety.cpp} & concurrency helpers & поведение на connection pool, thread-local достъп и transaction guard механизма \\
\path{test_task.cpp} & coroutine task abstraction & lifecycle на \code{Task<T>} и basic await/coroutine semantics \\
\path{test_thread_pool.cpp} & worker scheduling & posting, scheduling и worker resource coordination \\
\path{test_io_context.cpp} & async orchestration & event-loop и continuation plumbing за async execution слоя \\
\path{test_async_connector.cpp} & async connector facade & bridge логика между \code{async\_db<DB>} и coroutine-aware connector hooks \\
\path{test_wire_protocol.cpp} & experimental execution layer & ниско ниво на compile-time/wire-oriented механизми \\
\path{test_cpp26_reflection.cpp} & future language integration & early path за reflection-базирана автоматизация \\
\bottomrule
\end{tabularx}
\end{table}

Таблица~\ref{tab:verification_unit_core} е важна, защото показва, че unit слоят не се изчерпва с query parser-подобни проверки. Той покрива едновременно логическия модел, execution механиката, миграциите и coroutine-aware инфраструктурата, върху която стъпва асинхронната архитектура.

\begin{table}[H]
\centering
\small
\caption{Backend-oriented unit тестове и техният contract focus}
\label{tab:verification_unit_connectors}
\begin{tabularx}{\textwidth}{L{3.4cm}L{3.1cm}Z}
\toprule
\textbf{Файл} & \textbf{Backend фокус} & \textbf{Тип доказателство} \\
\midrule
\path{test_postgresql_connector.cpp} & \code{PostgreSQLDB} & \code{is\_connector} compliance, capability asserts и SQL parameter rendering \\
\path{test_mysql_connector.cpp} & \code{MySQLDB} & joins/transactions contract, dialect-specific rendering и bind expectations \\
\path{test_mongodb_connector.cpp} & \code{MongoDB} & negative capability validation и BSON-like filter rendering \\
\path{test_redis_connector.cpp} & \code{RedisDB} & absence of joins/transactions/aggregation и key-value command materialization \\
\path{test_cassandra_connector.cpp} & \code{CassandraDB} & capability ограничения и CQL rendering корелати \\
\path{test_neo4j_connector.cpp} & \code{Neo4jDB} & transaction capability и graph-oriented query materialization \\
\bottomrule
\end{tabularx}
\end{table}

Тези backend-oriented unit тестове са критични, защото доказват не само позитивни възможности, но и честно декларирани ограничения. При multi-backend библиотека отрицателната верификация --- че даден backend \emph{не} поддържа дадена операция --- е също толкова важна, колкото и позитивната.

\begin{table}[H]
\centering
\small
\caption{Каталог на integration и live тестовете}
\label{tab:verification_integration_tests}
\begin{tabularx}{\textwidth}{L{3.2cm}L{3.0cm}Y}
\toprule
\textbf{Файл} & \textbf{Режим на изпълнение} & \textbf{Какво демонстрира} \\
\midrule
\path{test_mockdb.cpp} & винаги наличен integration baseline & query composition, parameter binding, prepared queries и CRUD сценарии без външен backend \\
\path{test_sqlite.cpp} & локален integration тест & реален релационен execution path с capability asserts, joins, transactions и result hydration \\
\path{test_postgresql_live.cpp} & условен live тест & реална връзка, table lifecycle и CRUD изпълнение върху PostgreSQL \\
\path{test_mysql_live.cpp} & условен live тест & реален MySQL driver path, insert/update/delete и select filtering \\
\path{test_mongodb_live.cpp} & условен live тест & изпълнение на document-oriented сценарии срещу жив backend \\
\path{test_redis_live.cpp} & условен live тест & key-value и hash command materialization върху Redis среда \\
\path{test_cassandra_live.cpp} & условен live тест & wide-column execution и backend-specific CQL сценарии \\
\path{test_neo4j_live.cpp} & условен live тест & graph query execution през контролирана Docker среда \\
\bottomrule
\end{tabularx}
\end{table}

Таблица~\ref{tab:verification_integration_tests} показва, че repository-backed доказателството е разслоено. Не всички backend-и имат еднакъв operational cost за live verification, но хранилището съдържа ясна стратегия как и кога подобни тестове се активират. Това е особено важно за тезата, защото избягва смесването между mock correctness и реално driver-backed поведение.

\subsection{Карта на зрелостта според доказателствената опора}

Самият capability профил не е достатъчен, ако липсва evidence за реално поведение. Затова е полезно backend-ите да се разгледат и като стълбица на зрелост, определена не само от поддържаните операции, но и от наличната тестова опора. Тук зрелостта означава комбинация от contract correctness, execution evidence и operational реализъм.

\begin{table}[H]
\centering
\small
\caption{Backend maturity карта според наличната доказателствена опора}
\label{tab:verification_backend_maturity_map}
\begin{tabularx}{\textwidth}{L{2.4cm}L{2.2cm}L{2.9cm}Y}
\toprule
\textbf{Backend} & \textbf{Ниво на зрелост} & \textbf{Основна доказателствена опора} & \textbf{Интерпретация} \\
\midrule
\code{MockDB} & ниво 3 & \path{test_mockdb.cpp} & backend-independent execution baseline за query contract-а и prepared query механиката \\
SQLite & ниво 3 към 4 & \path{test_sqlite.cpp} & локален реален execution path с богата query evidence и capability asserts \\
PostgreSQL & ниво 3 & \path{test_postgresql_connector.cpp}, \path{test_postgresql_live.cpp} & силен contract плюс условна live backend проверка \\
MySQL & ниво 3 & \path{test_mysql_connector.cpp}, \path{test_mysql_live.cpp} & реален operational path с релационен профил и dialect-specific binding \\
MongoDB & ниво 3 & \path{test_mongodb_connector.cpp}, \path{test_mongodb_live.cpp} & document-oriented execution evidence без SQL-style capabilities \\
Redis & ниво 3 & \path{test_redis_connector.cpp}, \path{test_redis_live.cpp} & key-value execution evidence с ограничен, но честен contract профил \\
Neo4j & ниво 3 & \path{test_neo4j_connector.cpp}, \path{test_neo4j_live.cpp} & graph execution path с Docker-gated live verification \\
Cassandra & ниво 3 & \path{test_cassandra_connector.cpp}, \path{test_cassandra_live.cpp} & wide-column path с фокус върху формална пригодност и live coverage \\
\bottomrule
\end{tabularx}
\end{table}

Картата в Таблица~\ref{tab:verification_backend_maturity_map} е важна, защото предотвратява прекалено бинарно четене на backend статуса. Конекторът не е просто наличен или липсващ; той може да бъде contract-correct, execution-validated, partially live-verified или schema-aware богат. Именно този по-нюансиран прочит прави верификационната глава по-научно убедителна и по-близка до реалния статус на хранилището.

\subsection{Методология на benchmark оценката}

Емпиричната оценка е реализирана чрез Google Benchmark~\cite{google_benchmark_user_guide} в \path{benchmarks/bench_async.cpp}. Тя не цели да симулира пълния жизнен цикъл на конкретна база данни, а да изолира архитектурния въпрос дали coroutine-базираното изпълнение освобождава worker нишките по-ефективно от синхронното блокиране, когато между две кратки CPU фази има waiting интервал.

Benchmark наборът е разделен на две групи. Първата група измерва throughput на цели партиди от задачи. Синхронният baseline изпълнява CPU работа, след което worker нишката блокира за предварително зададен интервал. Асинхронният вариант използва \code{SimulatedAsyncIO} awaitable, който спира корутината, планира delayed resume чрез timer queue и връща продължаването обратно към \code{ThreadPool}. Втората група измерва изолирани overhead-и: създаване и изчакване на \code{Task<int>}, \code{ThreadPool::post()} round-trip и вложени \code{co\_await} вериги.

Методологично е важно, че throughput benchmark-ите използват wall-clock измерване, а не чисто CPU време. Причината е изрично отбелязана и в самия benchmark код: при корутинни сценарии с waiting фаза calibration по CPU време на macOS може да даде подвеждащо поведение, защото suspend-натата задача почти не консумира CPU. Затова съпоставката се извършва чрез \code{UseRealTime()} и фиксиран брой итерации.

В настоящата глава се използват release резултатите от локалния \path{build-rel} build tree. Те не трябва да се интерпретират като абсолютни production числа за всички възможни машини и драйвери. Тяхната роля е сравнителна: всички варианти се изпълняват върху една и съща машина, с една и съща параметрична мрежа, така че да се изведе относителният ефект от различния execution модел.

\begin{table}[H]
\centering
\small
\caption{Семейства benchmark измервания и тяхната цел}
\label{tab:benchmark_families}
\begin{tabularx}{\textwidth}{L{3.4cm}L{5.2cm}Y}
\toprule
\textbf{Benchmark семейство} & \textbf{Какво измерва} & \textbf{Основни параметри} \\
\midrule
\code{BM\_SyncThroughput} & throughput на блокиращ baseline, при който worker нишката остава заета по време на waiting фазата & брой нишки, брой задачи, waiting латентност, CPU iterations \\
\code{BM\_AsyncThroughput} & throughput на coroutine вариант със suspend/resume и delayed resume чрез timer queue & брой нишки, брой задачи, waiting латентност, CPU iterations \\
\code{BM\_AsyncYieldThroughput} & цена на coroutine interleave без изкуствена waiting латентност & брой нишки, брой задачи, CPU iterations \\
\code{BM\_TaskCreateAndWait} & минимална цена на създаване и завършване на тривиална задача & единична корутина без външно scheduling \\
\code{BM\_ThreadPoolPost} & цена на \code{post()} и promise/future round-trip към worker нишка & размер на thread pool-а \\
\code{BM\_NestedCoAwait} & overhead на вложено \code{co\_await} продължаване при различна дълбочина & дълбочина на coroutine веригата \\
\bottomrule
\end{tabularx}
\end{table}

\subsection{Основни throughput резултати}

Таблица~\ref{tab:throughput_results_release} обобщава представителните release-mode резултати за най-информативните конфигурации. Показани са двойки, за които има директно сравнение между синхронния и асинхронния вариант при едни и същи параметри.

\begin{table}[H]
\centering
\small
\caption{Representative release-mode throughput results for synchronous and asynchronous execution}
\label{tab:throughput_results_release}
\begin{tabularx}{\textwidth}{L{1.3cm}L{1.5cm}L{1.9cm}L{2.4cm}L{2.4cm}L{1.7cm}}
\toprule
\textbf{Нишки} & \textbf{Задачи} & \textbf{Waiting латентност} & \textbf{Синхронен throughput} & \textbf{Асинхронен throughput} & \textbf{Ускорение} \\
\midrule
4 & 100 & 10~$\mu$s & 194.6k ops/s & 658.5k ops/s & 3.38$\times$ \\
4 & 100 & 100~$\mu$s & 29.8k ops/s & 458.2k ops/s & 15.35$\times$ \\
4 & 100 & 1000~$\mu$s & 3.22k ops/s & 75.28k ops/s & 23.36$\times$ \\
4 & 1000 & 100~$\mu$s & 30.08k ops/s & 670.5k ops/s & 22.29$\times$ \\
8 & 100 & 100~$\mu$s & 57.13k ops/s & 258.5k ops/s & 4.53$\times$ \\
\bottomrule
\end{tabularx}
\end{table}

Най-важният извод е, че предимството на корутинния модел нараства не когато локалната CPU работа нарасне, а когато waiting фазата започне да доминира над нея. При 4 worker нишки и 100 задачи разликата между 10 и 1000 микросекунди ($\mu$s) waiting латентност увеличава относителното предимство на асинхронния вариант от 3.38$\times$ до 23.36$\times$. Това е именно очакваното поведение: при синхронния baseline нишките се изчерпват от блокиране, докато при корутинния модел те се освобождават за друга работа.

Вторият важен извод е, че ефектът не изчезва при по-големи партиди. При 4 нишки, 1000 задачи и 100~$\mu$s waiting латентност асинхронният вариант достига 670.5k операции в секунда срещу 30.08k при синхронния baseline. Това показва, че архитектурното предимство не е ограничено до „демо“ случай с малък брой задачи, а се запазва и когато конкуренцията за worker ресурси нарасне.

\begin{figure}[H]
\centering
\begin{tikzpicture}[x=1.95cm,y=0.24cm]
    \draw[->, thick] (0,0) -- (0,26) node[above] {\small ускорение async/sync};
    \draw[->, thick] (0,0) -- (6.4,0);

    \foreach \y in {5,10,15,20,25}
    {
        \draw[gray!35] (0,\y) -- (6.2,\y);
        \node[left] at (0,\y) {\scriptsize \y};
    }

    \foreach \x/\h/\labeltext/\valuetext in {
        0.7/3.38/{4 нишки\\100 задачи\\10~$\mu$s}/3.38,
        1.9/15.35/{4 нишки\\100 задачи\\100~$\mu$s}/15.35,
        3.1/23.36/{4 нишки\\100 задачи\\1000~$\mu$s}/23.36,
        4.3/22.29/{4 нишки\\1000 задачи\\100~$\mu$s}/22.29,
        5.5/4.53/{8 нишки\\100 задачи\\100~$\mu$s}/4.53}
    {
        \fill[green!28, draw=black] (\x-0.27,0) rectangle (\x+0.27,\h);
        \node[font=\scriptsize, align=center, above] at (\x,\h) {\valuetext$\times$};
        \node[font=\scriptsize, align=center, below=7pt] at (\x,0) {\labeltext};
    }
\end{tikzpicture}
\caption{Относително ускорение на асинхронния throughput спрямо синхронния baseline при representative I/O-bound конфигурации}
\label{fig:async_speedup_cases}
\end{figure}

Фигура~\ref{fig:async_speedup_cases} визуализира същата тенденция в относителна форма. Дори когато броят на worker нишките се удвои до 8, асинхронният вариант остава приблизително 4.5$\times$ по-бърз при 100~$\mu$s waiting латентност. Това е важен practically relevant резултат: повече нишки намаляват част от натиска върху синхронния baseline, но не елиминират структурното предимство на suspend/resume execution модела.

\subsection{Нулева waiting латентност (0ns) и coroutine scheduling overhead}

При нулева artificial waiting латентност сравнението престава да измерва освобождаването на worker нишките от I/O изчакване и започва да показва почти чистата цена на suspend/resume механиката. Поради това най-подходящата repository-backed съпоставка комбинира \code{BM\_SyncThroughput} с нулева waiting стойност и \code{BM\_AsyncYieldThroughput}, който заменя timer-базирания wait с еднократно \code{PoolYield} прехвърляне обратно към \code{ThreadPool}.

\begin{table}[H]
\centering
\small
\caption{Representative release-mode резултати при нулева artificial waiting латентност}
\label{tab:zero_latency_yield_results}
\begin{tabularx}{\textwidth}{L{1.4cm}L{1.5cm}L{2.8cm}L{2.9cm}Y}
\toprule
\textbf{Нишки} & \textbf{Задачи} & \shortstack{\textbf{Синхронен}\\\textbf{throughput}} & \shortstack{\textbf{Async yield}\\\textbf{throughput}} & \shortstack{\textbf{Отношение}\\\textbf{async/sync}} \\
\midrule
4 & 100 & 841.0k ops/s & 989.4k ops/s & 1.18$\times$ \\
4 & 1000 & 1.232M ops/s & 1.184M ops/s & 0.96$\times$ \\
8 & 100 & 810.9k ops/s & 514.4k ops/s & 0.63$\times$ \\
\bottomrule
\end{tabularx}
\end{table}

\begin{figure}[H]
\centering
\begin{tikzpicture}[x=1.65cm,y=4.0cm]
    \draw[->, thick] (0,0) -- (0,1.40) node[above] {\small throughput (M ops/s)};
    \draw[->, thick] (0,0) -- (5.8,0);

    \foreach \y/\label in {0.25/0.25,0.50/0.50,0.75/0.75,1.00/1.00,1.25/1.25}
    {
        \draw[gray!35] (0,\y) -- (5.6,\y);
        \node[left] at (0,\y) {\scriptsize \label};
    }

    \fill[blue!25, draw=black] (0.75,0) rectangle (1.05,0.841);
    \fill[green!28, draw=black] (1.15,0) rectangle (1.45,0.989);
    \node[font=\scriptsize, align=center, below=7pt] at (1.10,0) {4 нишки\\100 задачи};

    \fill[blue!25, draw=black] (2.55,0) rectangle (2.85,1.232);
    \fill[green!28, draw=black] (2.95,0) rectangle (3.25,1.184);
    \node[font=\scriptsize, align=center, below=7pt] at (2.90,0) {4 нишки\\1000 задачи};

    \fill[blue!25, draw=black] (4.35,0) rectangle (4.65,0.811);
    \fill[green!28, draw=black] (4.75,0) rectangle (5.05,0.514);
    \node[font=\scriptsize, align=center, below=7pt] at (4.70,0) {8 нишки\\100 задачи};

    \fill[blue!25, draw=black] (3.15,1.31) rectangle (3.33,1.37);
    \node[font=\scriptsize, anchor=west] at (3.37,1.34) {sync 0 wait};
    \fill[green!28, draw=black] (4.35,1.31) rectangle (4.53,1.37);
    \node[font=\scriptsize, anchor=west] at (4.57,1.34) {async yield};
\end{tikzpicture}
\caption{Throughput съпоставка при нулева artificial waiting латентност}
\label{fig:zero_latency_yield_throughput}
\end{figure}

Таблица~\ref{tab:zero_latency_yield_results} и Фигура~\ref{fig:zero_latency_yield_throughput} показват, че при липса на waiting фаза корутинният слой вече не получава структурния бонус, наблюдаван в I/O-bound режима. При 4 нишки и 100 задачи yield-only вариантът все още е приблизително 1.18$\times$ по-бърз, но при 4 нишки и 1000 задачи той е практически паритетен със синхронния baseline (0.96$\times$), а при 8 нишки и 100 задачи пада до 0.63$\times$. Това е полезен коректив към силните I/O-bound резултати: без външен wait coroutine machinery не „създава“ throughput, а добавя малка, но измерима scheduling цена, която може да бъде компенсирана само частично.

\subsection{Microbenchmark интерпретация на overhead-а}

Изолираните overhead измервания показват, че корутинната инфраструктура сама по себе си не е доминиращият разход в I/O-bound сценариите. В release build \code{BM\_TaskCreateAndWait} измерва приблизително 71.6~ns за създаване и завършване на тривиална задача. \code{BM\_ThreadPoolPost} остава около 5.0~$\mu$s независимо дали pool-ът има 1, 2, 4 или 8 worker нишки. \code{BM\_NestedCoAwait} достига около 324~ns при дълбочина 10 и около 2.375~$\mu$s дори при дълбочина 100.

Тези числа не трябва да се четат като пряка обещана цена на всяка реална заявка, но са важни с друго. Те показват, че overhead-ът на coroutine machinery е порядъци по-малък от 100--1000~$\mu$s waiting фази, при които асинхронният throughput печели най-много. С други думи, observed speedup-ът не е артефакт на измервателния инструмент, а следствие от това, че worker нишките престават да бъдат заключени в блокиращо изчакване.

\subsection{Ограничения на емпиричната оценка}

Емпиричната оценка съзнателно използва симулирана waiting фаза, а не пълно end-to-end изпълнение срещу конкретен мрежов сървър за база данни. Това е ограничение, ако целта е да се предскаже абсолютен throughput за реална PostgreSQL или Cassandra инсталация. Но то е и методологично предимство, ако целта е да се изолира execution моделът от влиянието на SQL planner, мрежовите колебания, кеширане в драйвера и storage engine спецификите.

Следователно настоящите benchmark резултати не доказват, че „асинхронното е винаги по-бързо“. Те доказват по-тясна, но по-научно защитима теза: когато архитектурата работи в I/O-bound режим и waiting фазата доминира над локалната CPU работа, coroutine-базираният модел позволява съществено по-добро използване на worker ресурсите от синхронния baseline. Това е точно твърдението, което асинхронният слой на библиотеката трябва да обоснове.

В комбинация с тестовата верификация този резултат затваря двата основни въпроса на дипломната работа. От една страна библиотеката показва формална и тестово подкрепена correctness дисциплина. От друга страна асинхронният execution модел показва измерима practically relevant полза именно в сценария, за който е проектиран. Така емпиричната оценка не стои встрани от архитектурата, а потвърждава нейния централен инженерeн и научен смисъл.
